\chapter{مثال من الواقع: معاملات الشركات المحدودة المسؤولية}
\label{ch:companies}

\section{القصة الحقيقية التي سنطبّقها}
\label{sec:co-story}
\textbf{لماذا؟} أتقنت الدورة مع فريق حماية المستهلك. الآن نطبّق الأدوات نفسها على معاملة تعالجها الوزارة فعلاً كل يوم: طلب تصديق النظام الأساسي لشركة محدودة المسؤولية، ثم طلب تسجيل تنازل عن حصص فيها.

في الواقع تُقدَّم طلبات تأسيس الشركات إلى «مديرية الشركات»، وهي تعمل في مبنى «الإدارة العامة للتجارة الداخلية وحماية المستهلك» بدمشق، وفيها «دائرة الشركات». تتحقق الدائرة من الوثائق وفق قانون الشركات الصادر بالمرسوم التشريعي رقم 29 لعام 2011، ثم يصدر قرار بالصيغة: «يُصدّق النظام الأساسي لشركة … المحدودة المسؤولية»، ويُنشر في الجريدة الرسمية. بعد ذلك تستكمل الشركة إشهارها في السجل التجاري.

\begin{table}[H]
\centering
\caption{المعاملتان في الواقع وفي كيان}
\label{tab:co-reality}
\begin{tabularx}{\linewidth}{R{3.4cm} Y Y}
\rcHeaderRow
\thh{المرحلة} & \thh{ما يحدث في الواقع} & \thh{ما يقابله في المراسلات}\\
تقديم الطلب & يقدّم وكيل المؤسسين الطلب مع النظام الأساسي وصور الهويات والوكالة المصدّقة. & وارد ورقي يُحفظ بانتظار القيد.\\
تسجيل الطلب & يُسجَّل في ديوان المديرية ويأخذ رقماً. & «قيد وتوزيع» على رئيسة دائرة الشركات للإجراء.\\
التدقيق & يدقق موظف الدائرة الوثائق ويبيّن النواقص. & إحالة إلى المدقق مع مهلة، واستفسار عند الشك.\\
النواقص & يُبلَّغ المراجع بما يلزم استكماله. & صادر «إشعار باستكمال النواقص» مرتبط بالوارد.\\
الاستكمال & يعود الوكيل بالوثائق الناقصة. & وارد جديد «رداً على كتابنا رقم».\\
التصديق & يصدر قرار تصديق النظام الأساسي. & صادر من نوع «قرار تصديق» يوقّعه مدير الشركات.\\
التنازل عن حصص & يُوقَّع التنازل أمام الكاتب بالعدل أو موظف المديرية وفق نموذج الوزارة، ثم يُقيَّد في سجل الشركاء. & وارد يُحال إلى أمينة سجل الشركاء، ثم صادر «إشعار بتسجيل التنازل».\\
\end{tabularx}
\end{table}

\begin{rcnote}[ما الحقيقي وما التدريبي؟]
أسماء الوزارة والإدارة العامة والمديرية والدائرة، وتسلسل المعاملة، وصيغة القرار، وشرط تسجيل التنازل في سجل الشركاء: كلها مأخوذة من الواقع. أسماء الموظفين والشركة والشركاء والأرقام والمبالغ تدريبية، وكل منها يحمل كلمة «تجريبي» أو «تدريب». لا تستخدم المبالغ هنا مرجعاً قانونياً؛ الحدود والرسوم النافذة تؤخذ من المديرية.
\end{rcnote}
\textbf{تحقق:} تستطيع أن تقول بجملة واحدة: الوارد طلب، والتدقيق إحالة، والنواقص صادر، والتصديق قرار موقّع، والتنازل لا يكتمل قبل قيده في سجل الشركاء.

\section{أضف وحدات مديرية الشركات في الموارد البشرية}
\label{sec:co-nodes}
\textbf{لماذا؟} لديك الآن «الإدارة العامة للتجارة الداخلية وحماية المستهلك» وتحتها فرع حماية المستهلك فقط. مديرية الشركات فرع شقيق تحت الإدارة العامة نفسها. نضيفها دون لمس وحدات حماية المستهلك.
\begin{enumerate}
\item بحساب المسؤول افتح «الموارد البشرية» ← «الوحدات التنظيمية» ← «إضافة».
\item أنشئ «مديرية الشركات» بالرمز \en{DIT-COMP} والنوع «مديرية»، واختر في «يتبع ل» «الإدارة العامة للتجارة الداخلية وحماية المستهلك»، ثم «إنشاء».
\item أنشئ «دائرة الشركات» بالرمز \en{DIT-COMP-REG} والنوع «دائرة»، واختر «مديرية الشركات» في «يتبع ل».
\item أنشئ «ديوان مديرية الشركات» بالرمز \en{DIT-COMP-DIWAN} والنوع «دائرة»، تابعاً أيضاً لـ«مديرية الشركات». سنربط به مكتب الديوان لاحقاً.
\end{enumerate}
\begin{table}[H]
\centering
\caption{الوحدات الجديدة}
\label{tab:co-nodes}
\begin{tabularx}{\linewidth}{Y R{3.2cm} Y}
\rcHeaderRow
\thh{الاسم ونوعه} & \thh{الرمز} & \thh{يتبع ل}\\
مديرية الشركات؛ مديرية & \en{DIT-COMP} & الإدارة العامة للتجارة الداخلية وحماية المستهلك\\
دائرة الشركات؛ دائرة & \en{DIT-COMP-REG} & مديرية الشركات\\
ديوان مديرية الشركات؛ دائرة & \en{DIT-COMP-DIWAN} & مديرية الشركات\\
\end{tabularx}
\end{table}
\textbf{تحقق:} في شجرة الوحدات ترى تحت الإدارة العامة فرعين: «مديرية حماية المستهلك وسلامة الغذاء» و«مديرية الشركات». إذا ظهرت مديرية الشركات في المستوى الأعلى، فقد نسيت «يتبع ل».

\section{أنشئ أوصاف الوظائف والمناصب}
\label{sec:co-positions}
\textbf{لماذا؟} المراسلات تُرسل من منصب، والقرار يوقّعه منصب. نحتاج خمسة كراسي: من يقود المديرية، ومن يرأس الدائرة، ومن يدقق الملفات، ومن يمسك سجل الشركاء، ومن يقيّد في الديوان.
\begin{enumerate}
\item افتح «وصف الوظيفة» ← «إضافة»، وأنشئ الأوصاف الخمسة بأسمائها ورموزها من الجدول.
\item افتح «المناصب» ← «إضافة». أنشئ المناصب من الأعلى إلى الأسفل حتى تجد المدير عند اختيار «المنصب الأعلى».
\item لكل منصب اختر الوصف والوحدة والمنصب الأعلى من الجدول، ثم «إنشاء».
\end{enumerate}
\begin{table}[H]
\centering
\caption{المناصب الخمسة}
\label{tab:co-positions}
\begin{tabularx}{\linewidth}{Y R{3.4cm} Y}
\rcHeaderRow
\thh{المنصب ورمزه} & \thh{الوصف} & \thh{الوحدة؛ المنصب الأعلى}\\
مدير الشركات؛ \en{POS-COMP-DIR} & \en{JD-COMP-DIR} & مديرية الشركات؛ فارغ، أو مدير الإدارة العامة إن وُجد منصبه.\\
رئيس دائرة الشركات؛ \en{POS-COMP-REG-H} & \en{JD-COMP-REG-H} & دائرة الشركات؛ مدير الشركات.\\
مدقق ملفات الشركات — 01؛ \en{POS-COMP-AUDITOR-01} & \en{JD-COMP-AUDITOR} & دائرة الشركات؛ رئيس دائرة الشركات.\\
أمين سجل الشركاء؛ \en{POS-COMP-PARTNERS-REG} & \en{JD-COMP-PARTNERS} & دائرة الشركات؛ رئيس دائرة الشركات.\\
رئيس ديوان مديرية الشركات؛ \en{POS-COMP-DIWAN-H} & \en{JD-COMP-DIWAN} & ديوان مديرية الشركات؛ مدير الشركات.\\
\end{tabularx}
\end{table}
\textbf{تحقق:} افتح «رئيس دائرة الشركات» وتأكد أن منصبه الأعلى «مدير الشركات». هذه التبعية هي التي يقترح منها النظام مسار الاعتماد لاحقاً.

\section{أضف الأعضاء وعيّن كل شخص في كرسيه}
\label{sec:co-members}
\textbf{لماذا؟} بلا تعيين فعّال لن يظهر المنصب في «أعمل بصفتي»، ولن يستطيع صاحبه الحفظ أو الإرسال. كل عضو يحتاج بريد دخول خاصاً به.
\begin{enumerate}
\item افتح «الأعضاء» ← «إضافة». أدخل الاسم ورقم الموظف وبريداً تدريبياً مختلفاً لكل عضو، وفعّل «التسجيل»، ثم «إنشاء».
\item افتح «التعيينات» ← «إضافة». اختر العضو والمنصب من الجدول، وتاريخ بدء اليوم أو قبله، والحالة «نشط»، ثم «إنشاء».
\item كرر للأعضاء الخمسة. لا تعيّن أحمد أو نور أو ريم هنا؛ يبقى فريق حماية المستهلك كما هو.
\end{enumerate}
\begin{table}[H]
\centering
\caption{الأعضاء الجدد ومهامهم في القصة}
\label{tab:co-members}
\begin{tabularx}{\linewidth}{R{3.6cm} Y Y}
\rcHeaderRow
\thh{العضو ورقمه} & \thh{المنصب} & \thh{ما يفعله في المعاملة}\\
سامر العلي — تجريبي؛ \en{DEMO-COMP-001} & مدير الشركات & يوقّع قرار التصديق.\\
لينا درويش — تجريبي؛ \en{DEMO-COMP-002} & رئيس دائرة الشركات & تستلم الوارد وتحيله، تجيب الاستفسار، تراجع القرار، وتوقّع الإشعارات.\\
باسل يوسف — تجريبي؛ \en{DEMO-COMP-003} & مدقق ملفات الشركات — 01 & يدقق الوثائق، ويكتب إشعار النواقص ومشروع القرار.\\
هبة منصور — تجريبي؛ \en{DEMO-COMP-004} & أمين سجل الشركاء & تقيّد التنازل في سجل الشركاء وتكتب إشعاره، وترفع سجلات الأرشفة.\\
خالد عيسى — تجريبي؛ \en{DEMO-COMP-005} & رئيس ديوان مديرية الشركات & يقيّد الوارد، ويرسل الصادر ويسجّل تسليمه.\\
\end{tabularx}
\end{table}
\begin{rcnote}[اختصار للتدريب على حساب واحد]
إن لم تتوفر خمسة حسابات، عيّن حسابك التدريبي في أكثر من منصب، وبدّل «أعمل بصفتي» قبل كل إجراء. هذا يكفي لفهم المسار، لكنه لا يختبر الصلاحيات كما يختبرها شخص حقيقي؛ الاختبار الكامل بخمسة حسابات.
\end{rcnote}
\textbf{تحقق:} ادخل بحساب باسل وافتح «المراسلات الرسمية». يجب أن يظهر «مدقق ملفات الشركات — 01» في «أعمل بصفتي»، وألا تظهر رسالة «لا يمكنك حفظ المراسلات أو إرسالها بعد».

\section{امنح أدوار المراسلات لمن يحتاجها فقط}
\label{sec:co-roles}
\textbf{لماذا؟} شغل المنصب يكفي للكتابة والاعتماد والإحالة. أما القيد في الديوان والأرشفة فلهما صلاحيات خاصة. نمنح كل صلاحية لمنصب واحد واضح.
\begin{enumerate}
\item بحساب المسؤول افتح «إعدادات المراسلات» ← «الأدوار والصلاحيات» ← «إدارة الأدوار» ← «إضافة الدور».
\item أنشئ الأدوار الثلاثة من الجدول. في كل دور افتح «الصلاحيات» وفعّل الصلاحية المذكورة فقط ثم «حفظ»، ثم افتح «المسؤولين» وأسند المنصب.
\item لا تمنح «إدارة المراسلات» لأي عضو من الخمسة؛ تبقى عند مسؤول الإعداد.
\end{enumerate}
\begin{table}[H]
\centering
\caption{الأدوار الثلاثة}
\label{tab:co-roles}
\begin{tabularx}{\linewidth}{Y R{3.6cm} Y}
\rcHeaderRow
\thh{عنوان الدور} & \thh{الصلاحية} & \thh{يُسند إلى}\\
ديوان مديرية الشركات & كاتب الديوان & رئيس ديوان مديرية الشركات\\
أرشيف دائرة الشركات & أمين أرشيف المراسلات & أمين سجل الشركاء\\
اعتماد أرشيف الشركات & اعتماد الأرشيف & رئيس دائرة الشركات\\
\end{tabularx}
\end{table}
\textbf{تحقق:} يظهر كل منصب في «الكيانات المعينة» لدوره. من رفع سجل الأرشفة (هبة) ليس من يعتمده (لينا).

\section{جهّز الديوان والأنواع والترويسة والجهات}
\label{sec:co-settings}
\textbf{لماذا؟} الوارد يحتاج مكتب ديوان وسلسلة أرقام، والقرار يحتاج نوعاً يفرض توقيع المدير، والكتاب يحتاج ترويسة رسمية، والوكيل يحتاج سجلاً في دليل الجهات.
\begin{enumerate}
\item «إعدادات المراسلات» ← «مكاتب الديوان»: اختر «ديوان مديرية الشركات» ثم «إضافة مكتب»، واجعله «فعّال».
\item «الترقيم»: بادئة الوارد \en{COMP-IN} وبادئة الصادر \en{COMP-OUT}، بأربع خانات، مع «السنة في الرقم». احفظ كل نمط.
\item «أنواع المراسلات» ← «إضافة نوع»: أنشئ النوعين في الجدول التالي ثم «حفظ النوع».
\item «الترويسات والأختام» ← «ترويسة جديدة»: اكتب الأسطر الأربعة من الجدول، والتذييل «نسخة تدريبية»، واجعلها متاحة لمديرية الشركات، ثم اربطها بالمديرية من «ترويسات الوحدات».
\item «الجهات الخارجية» ← «جهة جديدة»: أضف «وكيل مؤسسي شركة غيم الشام للتجارة العامة — تدريب». اترك البريد فارغاً؛ المسار ورقي كما في الواقع.
\end{enumerate}
\begin{table}[H]
\centering
\caption{النوعان والترويسة}
\label{tab:co-types}
\begin{tabularx}{\linewidth}{R{3.8cm} Y Y}
\rcHeaderRow
\thh{الحقل} & \thh{معاملة شركة محدودة المسؤولية} & \thh{قرار تصديق نظام أساسي}\\
المفتاح & \en{comp-llc-case} & \en{comp-llc-ratification}\\
البادئة & \en{COMP} & \en{COMP-DEC}\\
الأولوية والسرية & متوسط؛ عام & متوسط؛ عام\\
المهلة (أيام عمل) & 5 & 2\\
يتطلب مسار اعتماد & مفعّل: المدقق يكتب، ورئيسة الدائرة توقّع. & مفعّل: المدقق يكتب، ورئيسة الدائرة تراجع، والمدير يوقّع.\\
يمر عبر الديوان & مفعّل & مفعّل\\
اعتماد الأرشيف & مطلوب & مطلوب\\
\end{tabularx}
\end{table}
\begin{rcnote}[الترويسة كما في الكتب الرسمية]
السطر الأول «الجمهورية العربية السورية»، ثم «وزارة الاقتصاد والصناعة»، ثم «الإدارة العامة للتجارة الداخلية وحماية المستهلك»، ثم «مديرية الشركات». الرقم والتاريخ يملؤهما النظام عند الإصدار؛ لا تكتبهما يدوياً في النص.
\end{rcnote}
\textbf{تحقق:} يفتح خالد «ديوان الهيئة» ويرى مكتب مديرية الشركات، ويرى باسل النوعين في «مراسلة جديدة»، وتظهر الترويسة بأسطرها الأربعة في مسودة جديدة.

\section{قيّد طلب تصديق النظام الأساسي}
\label{sec:co-incoming}
\textbf{لماذا؟} هكذا تبدأ المعاملة في الواقع: وكيل المؤسسين يسلّم الملف الورقي إلى الديوان، فيُسجَّل ويُوجَّه إلى الدائرة المختصة.
\begin{enumerate}
\item جهّز ملف PDF تدريبياً بعنوان «طلب تصديق النظام الأساسي لشركة غيم الشام للتجارة العامة المحدودة المسؤولية — تدريب». ضع فيه: صفحة الطلب، والنظام الأساسي، وصوراً تدريبية للهويات، ووكالة خاصة غير مصدّقة. سنكتشف هذا النقص عمداً.
\item بحساب خالد افتح «ديوان الهيئة» ← «قيد وارد ورقي». ارفع الملف، واكتب الموضوع ورقم المرسل \en{GHAIM-01}، ثم «حفظ بانتظار القيد».
\item افتح البطاقة من «بانتظار القيد» ثم «قيد وتوزيع» بالقيم التالية، واضغط «قيد وتوزيع».
\end{enumerate}
\begin{table}[H]
\centering
\caption{بيانات القيد}
\label{tab:co-incoming}
\begin{tabularx}{\linewidth}{R{3.7cm} Y}
\rcHeaderRow
\thh{الحقل} & \thh{القيمة}\\
الجهة المرسلة & \req{} وكيل مؤسسي شركة غيم الشام للتجارة العامة — تدريب\\
الموضوع & \req{} طلب تصديق النظام الأساسي لشركة غيم الشام للتجارة العامة المحدودة المسؤولية — تدريب\\
النوع & معاملة شركة محدودة المسؤولية\\
يوزَّع إلى & \req{} رئيس دائرة الشركات، «للإجراء»\\
حاشية الديوان & \rec{} «للتدقيق وفق قانون الشركات وبيان الرأي»\\
\end{tabularx}
\end{table}
\textbf{تحقق:} للطلب رقم وارد يبدأ بـ\en{COMP-IN}، وهو منفصل عن \en{GHAIM-01}. يظهر عمل مفتوح في صندوق لينا.

\section{أحل الملف إلى المدقق مع مهلة}
\label{sec:co-refer}
\textbf{لماذا؟} رئيسة الدائرة لا تدقق كل ملف بنفسها؛ تحيله إلى المدقق وتتابع نتيجته.
\begin{enumerate}
\item بحساب لينا وبصفة «رئيس دائرة الشركات» افتح الوارد واضغط «إحالة».
\item اختر «مدقق ملفات الشركات — 01»، والغرض «للإجراء»، ومهلة يومي عمل.
\item اكتب الحاشية: «يرجى تدقيق النظام الأساسي والوثائق المرفقة وبيان النواقص إن وُجدت».
\item أكّد «إحالة» وافتح «المسار».
\end{enumerate}
\textbf{تحقق:} يرى باسل عملاً مفتوحاً بمهلة يومين، وترى لينا محطة الإحالة باسم منصبه.

\section{دقق الوثائق واطلب إفادة عند الشك}
\label{sec:co-audit}
\textbf{لماذا؟} القانون يطلب وثيقة تعيين المدير إذا لم يُعيَّن في النظام الأساسي. في ملفنا ورد اسم المدير داخل النظام، فهل نطلب الوثيقة؟ القرار عند رئيسة الدائرة، ويجب أن يُحفظ مع الملف.
\begin{enumerate}
\item بحساب باسل افتح الوارد وقارن الملف بقائمة التدقيق أدناه.
\item اضغط «استفسار»، واختر منصب لينا، واكتب: «ورد تعيين مدير الشركة في النظام الأساسي. هل نكتفي بذلك أم نطلب وثيقة تعيين مستقلة؟ وألاحظ أن الوكالة الخاصة غير مصدّقة».
\item بحساب لينا افتح «الاستفسارات» ثم «إفادة»: «نكتفي بالنظام الأساسي لتعيين المدير. تُطلب الوكالة الخاصة مصدّقة أصولاً».
\end{enumerate}
\begin{table}[H]
\centering
\caption{قائمة تدقيق ملف التأسيس}
\label{tab:co-checklist}
\begin{tabularx}{\linewidth}{R{4.2cm} Y}
\rcHeaderRow
\thh{الوثيقة} & \thh{النتيجة في ملفنا}\\
طلب التأسيس & موجود.\\
النظام الأساسي & موجود، ويتضمن اسم المدير ورأس المال والحصص.\\
صور الهويات أو جوازات السفر & موجودة لجميع الشركاء.\\
وثيقة تعيين المدير & غير لازمة؛ التعيين وارد في النظام الأساسي.\\
وكالة خاصة بالتأسيس مصدّقة & \req{} ناقصة: الوكالة غير مصدّقة.\\
\end{tabularx}
\end{table}
\textbf{تحقق:} سُجّل السؤال والإفادة داخل الوارد نفسه، وزال انتظار الإفادة. أصبح النقص الوحيد محدداً: الوكالة المصدّقة.

\section{أرسل إشعار النواقص بكتاب مرتبط}
\label{sec:co-deficiency}
\textbf{لماذا؟} المراجع يحتاج كتاباً رسمياً برقم يذكر ما ينقص. «رد» يربط الإشعار بالطلب، فيظهر كل شيء معاً عند عودة الوكيل.
\begin{enumerate}
\item بحساب باسل افتح الوارد واضغط «رد». اختر النوع «معاملة شركة محدودة المسؤولية» وترويسة المديرية.
\item الموضوع: «استكمال وثائق طلب تصديق النظام الأساسي لشركة غيم الشام — تدريب». الصق النص أدناه.
\item تأكد أن «إلى» هي جهة الوكيل، وأن المسار مرحلة واحدة: لينا بدور «توقيع». احفظ وأعد الفتح، ثم «إرسال للاعتماد».
\item توقّع لينا من «بانتظار اعتمادي» بـ«اعتماد وتوقيع». ثم يسجّل خالد من «بانتظار الإرسال» «تسجيل الإرسال» بالوسيلة الورقية بعد تسليم النسخة للوكيل.
\end{enumerate}
\begin{rcnote}[نص الإشعار]
السيد وكيل مؤسسي شركة غيم الشام للتجارة العامة المحدودة المسؤولية المحترم. إشارة إلى طلبكم المسجل لدينا برقم الوارد المذكور في المرجع، المتضمن تصديق النظام الأساسي للشركة، نعلمكم بضرورة تقديم الوكالة الخاصة بالتأسيس مصدّقة أصولاً، وذلك لمتابعة إجراءات التصديق. مع التقدير. رئيس دائرة الشركات.
\end{rcnote}
\textbf{تحقق:} للإشعار رقم \en{COMP-OUT}، ويظهر في الوارد رابط إليه، وفي «سجل الصادر» وسيلة الإرسال ونتيجتها. الإصدار ليس تسليماً؛ لا تسجّل التسليم قبل حدوثه.

\section{استقبل الوكالة المصدّقة واربطها بالإشعار}
\label{sec:co-complete}
\textbf{لماذا؟} عندما يعود الوكيل بالوثيقة، يجب أن يعرف الجميع أنها جواب على إشعارنا، لا طلب جديد.
\begin{enumerate}
\item بحساب خالد قيّد وارداً ورقياً جديداً: «تقديم الوكالة الخاصة المصدّقة — شركة غيم الشام — تدريب»، برقم مرسل \en{GHAIM-02}.
\item في «رداً على كتابنا رقم» اختر رقم إشعار النواقص الفعلي، ووزّعه على «مدقق ملفات الشركات — 01» «للإجراء».
\item بحساب باسل افتح الوارد وتأكد من الوكالة. الملف الآن مكتمل.
\end{enumerate}
\textbf{تحقق:} يظهر الوارد الجديد مرتبطاً بالإشعار، والإشعار مرتبطاً بالطلب الأصلي: سلسلة واحدة من ثلاثة كتب.

\section{اكتب قرار التصديق وأرسله في مساره}
\label{sec:co-decision}
\textbf{لماذا؟} هذا هو الكتاب الذي ينتظره المؤسسون. يكتبه المدقق، وتراجعه رئيسة الدائرة، ويوقّعه المدير. في الواقع يصدر القرار باسم الوزير أو من فوّضه؛ في التدريب يوقّعه مدير الشركات بوصفه المفوَّض.
\begin{enumerate}
\item بحساب باسل افتح الوارد \en{GHAIM-02} واضغط «رد»، واختر النوع «قرار تصديق نظام أساسي».
\item الموضوع: «تصديق النظام الأساسي لشركة غيم الشام للتجارة العامة المحدودة المسؤولية — تدريب». الصق نص القرار أدناه، واكتب عمداً رأس المال «40,000,000» بدلاً من «50,000,000» الوارد في النظام الأساسي.
\item في «مسار الاعتماد» أضف: «رئيس دائرة الشركات» بدور «مراجعة»، ثم «مدير الشركات» بدور «توقيع». أرفق النظام الأساسي. احفظ وأعد الفتح ثم «إرسال للاعتماد».
\end{enumerate}
\begin{rcnote}[نص القرار]
إن مدير الشركات، بناءً على أحكام قانون الشركات الصادر بالمرسوم التشريعي رقم 29 لعام 2011، وعلى الطلب المقدم من وكيل المؤسسين والوثائق المرفقة به، يقرر ما يلي: المادة 1: يُصدّق النظام الأساسي لشركة غيم الشام للتجارة العامة المحدودة المسؤولية — تدريب، برأس مال قدره 50,000,000 ليرة سورية. المادة 2: على الشركة استكمال إجراءات الإشهار في السجل التجاري خلال المدة القانونية. المادة 3: يُبلَّغ هذا القرار من يلزم لتنفيذه.
\end{rcnote}
\textbf{تحقق:} يظهر القرار في «بانتظار اعتمادي» لدى لينا أولاً، لا لدى سامر.

\begin{figure}[H]
\centering
\begin{tikzpicture}[node distance=0.5cm]
\node[stepBox] (a) {\RL{خالد: قيد الطلب وتوزيعه}};
\node[stepBox,below=of a] (b) {\RL{لينا: إحالة إلى المدقق}};
\node[stepBox,below=of b] (c) {\RL{باسل: تدقيق وإشعار نواقص}};
\node[stepBox,below=of c] (d) {\RL{باسل: مشروع القرار}};
\node[stepBox,below=of d] (e) {\RL{لينا: مراجعة}};
\node[stepBox,below=of e] (f) {\RL{سامر: توقيع القرار}};
\node[stepBox,below=of f] (g) {\RL{خالد: إرسال وتسليم}};
\draw[flowsTo] (a) -- (b);
\draw[flowsTo] (b) -- (c);
\draw[flowsTo] (c) -- (d);
\draw[flowsTo] (d) -- (e);
\draw[flowsTo] (e) -- (f);
\draw[flowsTo] (f) -- (g);
\end{tikzpicture}
\rclegend
\caption{مسار طلب التصديق في مديرية الشركات}
\label{fig:co-route}
\shotnote{رسم توضيحي لتسلسل المسؤوليات، وليس لقطة من التطبيق. بين التدقيق ومشروع القرار يخرج إشعار النواقص ويعود الوكيل بالوثيقة. إذا أعادت لينا القرار للتعديل يرجع إلى باسل، ثم يبدأ المسار من مراجعتها من جديد.}
\end{figure}

\section{صحّح الخطأ بالإعادة ثم وقّع القرار}
\label{sec:co-sign}
\textbf{لماذا؟} خطأ في رقم رأس المال يجعل القرار مخالفاً للنظام الأساسي. الإعادة للتعديل تصلحه قبل التوقيع؛ بعد التوقيع يصبح النص ثابتاً.
\begin{enumerate}
\item بحساب لينا افتح القرار وقارنه بالنظام الأساسي المرفق. اضغط «إعادة للتعديل»: «رأس المال في المادة 1 لا يطابق النظام الأساسي؛ المطلوب 50,000,000 ليرة سورية».
\item بحساب باسل صحّح الرقم، واحفظ، وأعد الإرسال.
\item بحساب لينا اضغط «اعتماد». بحساب سامر وبصفة «مدير الشركات» افتح «بانتظار اعتمادي» ثم «اعتماد وتوقيع».
\item بحساب خالد افتح «بانتظار الإرسال»، وسجّل الإرسال الورقي بعد تسليم نسخة القرار للوكيل.
\end{enumerate}
\begin{rcnote}[بعد القرار، خارج المراسلات]
يُنشر قرار التصديق في الجريدة الرسمية، وتستكمل الشركة إشهارها في السجل التجاري خلال المدة القانونية. هذه خطوات الشركة والجهات المختصة؛ لا تنتظرها لإنجاز عمل الدائرة على الطلب.
\end{rcnote}
\textbf{تحقق:} للقرار رقم صادر وتوقيع سامر، ويحمل الرقم الصحيح لرأس المال. في «المسار» ترى الإعادة وسببها. الطلب الأصلي لم يعد عملاً مفتوحاً لدى لينا.

\section{افتح إضبارة الشركة وأرشف المعاملة}
\label{sec:co-dossier}
\textbf{لماذا؟} في الواقع لكل شركة ملف يجمع كل ما يخصها. ستعود إليه عند أول تنازل أو تعديل.
\begin{enumerate}
\item بحساب باسل افتح «الأضابير» ← «إضبارة جديدة». الوحدة «دائرة الشركات»، والعنوان «شركة غيم الشام للتجارة العامة المحدودة المسؤولية — تدريب»، ثم «فتح الإضبارة».
\item من كل كتاب في السلسلة اختر «إجراءات أخرى» ← «إضافة إلى إضبارة»: الطلب، والإشعار، والوكالة، والقرار.
\item «مشاركة» الإضبارة مع منصبي لينا وهبة. لا تغلق الإضبارة؛ حياة الشركة لم تنته.
\item بحساب هبة افتح «الأرشيف» ← «بانتظار المراجعة». راجع سجل القرار وأضف الوسمين «تصديق نظام أساسي» و«تدريب»، ثم «رفع للاعتماد». تعتمده لينا من «بانتظار الاعتماد» ← «اعتماد نهائي».
\end{enumerate}
\textbf{تحقق:} الكتب الأربعة في الإضبارة، وسجل أرشفة القرار «نهائي» برقم أرشيف. الإضبارة مفتوحة.

\section{طبّق المعاملة الثانية: تسجيل تنازل عن حصص}
\label{sec:co-transfer}
\textbf{لماذا؟} بعد التأسيس يقرر الشريك «عمر — تدريب» التنازل عن حصصه للشريك «مازن — تدريب». القانون يجعل التنازل أمام الكاتب بالعدل أو موظف المديرية وفق نموذج الوزارة، ولا يصبح المتنازَل له شريكاً قبل قيد الانتقال في سجل الشركاء.
\begin{enumerate}
\item بحساب خالد قيّد وارداً ورقياً: «طلب تسجيل تنازل عن حصص في شركة غيم الشام — تدريب» من الجهة نفسها، برقم مرسل \en{GHAIM-03}. أرفق عقد تنازل تدريبياً على نموذج الوزارة، موقّعاً أمام الكاتب بالعدل. وزّعه على لينا «للإجراء».
\item بحساب لينا أحله إلى «أمين سجل الشركاء» مع حاشية: «للتحقق من عقد التنازل وقيده في سجل الشركاء» ومهلة يومي عمل.
\item بحساب هبة افتح الوارد وتأكد من: النموذج المعتمد، وتوقيع الطرفين، وجهة التصديق، وعدد الحصص. ثم سجّل القيد في سجل الشركاء خارج المراسلات بحسب إجراءات الدائرة.
\item اضغط «رد» واختر «معاملة شركة محدودة المسؤولية». الموضوع: «إشعار بتسجيل التنازل عن حصص في سجل الشركاء — شركة غيم الشام — تدريب». النص: «إشارة إلى طلبكم المسجل لدينا، نعلمكم بأنه تم قيد التنازل عن الحصص المذكورة في العقد المرفق بطلبكم في سجل الشركاء أصولاً. مع التقدير».
\item المسار: لينا بدور «توقيع». بعد التوقيع يسجّل خالد الإرسال، وتضيف هبة الطلب والإشعار إلى إضبارة الشركة.
\end{enumerate}
\begin{rcnote}[جرّب هذا الخطأ الشائع]
لو أرسلت هبة الإشعار قبل القيد في سجل الشركاء، لأبلغت الشركة بشيء لم يحدث. القاعدة: القيد أولاً، ثم الإشعار. وإذا كان العقد على غير نموذج الوزارة، يكون الرد إشعار نواقص كما في المعاملة الأولى، لا إشعار تسجيل.
\end{rcnote}
\textbf{تحقق:} إضبارة الشركة تضم الآن ستة كتب في معاملتين، ولكل معاملة سلسلة مرتبطة من الطلب إلى الجواب.

\section{قائمة البطل: ما الذي يجب أن تستطيعه الآن}
\label{sec:co-hero}
\textbf{لماذا؟} تصبح مرجع فريقك عندما تشرح لماذا جرى كل شيء، لا أين الأزرار فقط. أجب عن هذه الأسئلة من التطبيق، دون الرجوع إلى الدليل.
\begin{table}[H]
\centering
\caption{أثبت إتقانك}
\label{tab:co-hero}
\begin{tabularx}{\linewidth}{Y Y}
\rcHeaderRow
\thh{السؤال} & \thh{أين تجد الجواب؟}\\
من يملك الطلب الآن، ومنذ متى؟ & «المسار» في الطلب الأصلي.\\
لماذا لم نطلب وثيقة تعيين المدير؟ & الإفادة المسجلة داخل الوارد.\\
لماذا عاد القرار إلى باسل؟ & سبب «إعادة للتعديل» في المسار.\\
هل استلم الوكيل القرار فعلاً؟ & «سجل الصادر»: وسيلة الإرسال ونتيجته، لا رقم القرار وحده.\\
ما كل ما يخص شركة غيم الشام؟ & إضبارتها المفتوحة.\\
متى صار مازن شريكاً؟ & عند قيد التنازل في سجل الشركاء، لا عند تقديم الطلب ولا عند توقيع الإشعار.\\
\end{tabularx}
\end{table}
\begin{itemize}
\item كرر المعاملة الأولى بشركة جديدة دون نواقص، وقِس الوقت من القيد إلى التسليم.
\item فعّل من «التحويل الآلي» قاعدة توجّه وارد مديرية الشركات إلى «رئيس دائرة الشركات»، ثم اختبرها بطلب جديد.
\item اشرح لزميل الفرق بين «إعادة للتعديل» و«رفض»، وبين «للإجراء» و«للاطلاع»، من أمثلة هذا الفصل.
\end{itemize}
\textbf{تحقق:} أجبت عن الأسئلة الستة من التطبيق، ونفّذت تمريناً واحداً على الأقل دون الدليل.
